\section{Calidad de los datos de simulación: cómo abordar la sim-to-real gap}

La simulación es la capa intermedia más importante de la pirámide de datos para robots. Permite generar escenarios a escala, controlar variables, reproducir fallas y construir evaluaciones. Sin embargo, los datos de simulación solo tienen valor cuando la diferencia con el mundo real está bajo control. De lo contrario, el modelo aprende los atajos de la simulación y no las capacidades que necesita en la realidad.

\subsection{Las cuatro dimensiones de la calidad de la simulación}

\noindent\begin{tabularx}{\textwidth}{p{0.22\textwidth}p{0.34\textwidth}X}
\toprule
Dimensión & Requisito & Manifestación de la falla \\
\midrule
Coherencia física & Contacto, fricción, gravedad y dinámica creíbles & Las acciones fallan en la realidad. \\
Realismo visual & Materiales, iluminación, oclusiones y ruido de cámara razonables & El modelo de percepción se sobreajusta a las imágenes de la simulación. \\
Diversidad de tareas & Cubrir la cola larga, las anomalías y las combinaciones complejas & Solo resuelve escenarios estándar. \\
Transferibilidad de la evaluación & La evaluación en simulación predice el desempeño real & Las tablas de clasificación en simulación no corresponden con el despliegue real. \\
\bottomrule
\end{tabularx}

\begin{warningbox}{Más realismo no siempre es mejor simulación}
Perseguir en exceso el realismo visual puede tener un costo alto y un beneficio bajo. La simulación debe servir a los objetivos de entrenamiento y evaluación: hay que distinguir con claridad qué variables influyen en la política y cuáles solo afectan la apariencia.
\end{warningbox}

\subsection{Resumen de la sección}

La calidad de los datos de simulación no depende de «si se ve parecido o no», sino de si mejora el desempeño en tareas reales y ayuda a evaluar las capacidades reales.
